% !PRINTED 379
\weberpage{379}
%<<CONT
Es durchlaufe $\varepsilon$ die sämtlichen Teiler von $m$ (1 und $m$
eingeschlossen) und $e^2$ die sämtlichen quadratischen Teiler von
$m$; dann ist

\weq{10}{\overset{e^2}{\Sigma}\ \psi\!\left(D,\frac{m}{e^2}\right) = \overset{\varepsilon}{\Sigma}\ (D,\varepsilon).}

Ist diese Relation richtig für zwei Zahlen $m'$, $m''$, die keinen
gemeinsamen Teiler haben, so folgt sie für $m=m'm''$. Denn
haben $e'$, $\varepsilon'$ und $e''$, $\varepsilon''$ dieselbe Bedeutung für $m'$ und $m''$, wie
$e$, $\varepsilon$ für $m$, so sind $\varepsilon=\varepsilon'\varepsilon''$ die Teiler und $e^2=e'^2e''^2$ die
quadratischen Teiler von $m=m'm''$. Es ist aber nach (6) und
\S\,85, (8):

\[ \overset{e'^2}{\Sigma}\ \overset{e''^2}{\Sigma}\ \psi\!\left(D,\frac{m'}{e'^2}\right)\psi\!\left(D,\frac{m''}{e''^2}\right) = \overset{e^2}{\Sigma}\ \psi\!\left(D,\frac{m}{e^2}\right), \]

\[ \overset{\varepsilon'}{\Sigma}\ \overset{\varepsilon''}{\Sigma}\ (D,\varepsilon')(D,\varepsilon'') = \Sigma(D,\varepsilon). \]

Hiernach brauchen wir die Relation (10) nur noch unter der
Voraussetzung zu beweisen, daß $m=q^k$ eine Primzahlpotenz ist.

Unter dieser Voraussetzung ist aber:

\[
\begin{aligned}
\frac{n}{e^2}&=q^{k-2s}, & 0&\leqq s \leqq \frac{k}{2},\\
\varepsilon&=q^s, & s&=0,1,2,\dots,k.
\end{aligned}
\] % [sic] Gedruckt ist im Zähler n statt m; das zweite Relationszeichen ist eine abweichende Type.

Ist nun

1. $(D,q)=0$, so ist $\overset{s}{\Sigma}\ \psi(D,q^{k-2s}) = 1$,

2. $(D,q)=-1$, \text{\glqq}\ \text{\glqq}\ $\overset{s}{\Sigma}\ \psi(D,q^{k-2s}) = 1$ ($k$ gerade),
\hspace{4em}$= 0$ ($k$ ungerade),

3. $(D,q)=1$, \text{\glqq}\ \text{\glqq}\ $\overset{s}{\Sigma}\ \psi(D,q^{k-2s}) = k+1$,

und ebenso groß ergibt sich in den drei Fällen die Summe

\[ \underset{0,\,k}{\overset{s}{\Sigma}}\ (D,q^s); \]

denn im Falle 1. hat nur das dem $s=0$ entsprechende Glied
dieser Summe den Wert 1, die anderen verschwinden; im Falle 2.
haben die den geraden $s$ entsprechenden Glieder den Wert $+1$,
die anderen den Wert $-1$, und im Falle 3. haben alle
$k+1$-Glieder den Wert $+1$.

Damit ist also die Relation (10) allgemein bewiesen.
